The local \ac{sa} is structured around the optimized parameter set derived in Assignment~1. The main code structure is outlined in table~\ref{tab:code_02}.

\begin{table}[!ht]
\centering
\begin{tabular}{|l|}
\hline
read input: time series and catchment area (*.csv) \\ \hline
define optimized \ac{mp} set from Assignment~1 \\ \hline
define \ac{mp} boundaries \\ \hline
generate perturbation range: $x \in [-30\%, +30\%]$, 120 linearly spaced steps \\ \hline
\begin{tabular}[c]{@{}l@{}}
for each \ac{mp} $i$ and each perturbation $j$: \\
\quad -- compute perturbed value $x_{i,j}$, clipped to parameter bounds \\
\quad -- run HBV001a with the perturbed parameter set \\
\quad -- compute $\Delta \text{OFV}$ relative to the reference \ac{nse} \\
\quad -- record internal model outputs at $-30\%, -20\%, -10\%, +10\%, +20\%, +30\%$
\end{tabular} \\ \hline
document cases where perturbed values hit parameter boundaries \\ \hline
generate scatter plots: relative parameter change vs.\ relative \ac{ofv} change \\ \hline
rank \acp{mp} by sensitivity \\ \hline
\end{tabular}
\caption{Local \ac{sa}: main code structure}
\label{tab:code_02}
\end{table}

For each of the 18 \acp{mp}, the perturbation is applied as a relative change to the optimized value, $x_{i,j} = p_i \cdot (1 + j/100)$, where $j$ is drawn from 120 equally spaced steps between $-30\%$ and $+30\%$. The exception is \texttt{snw\_att}, whose optimized value is close to zero; here, an absolute perturbation of $\pm 0.25\,\si{K}$ is applied instead to avoid numerical issues with relative scaling. The perturbed value is always clipped to the defined parameter bounds before the model is run.

The sensitivity is quantified by the relative change in the \ac{ofv}:
\begin{equation}
    \Delta \text{OFV}_{i,j} \;=\; \frac{\text{OFV}_{i,j} - \text{OFV}_0}{\text{OFV}_0} \times 100\,\%
\end{equation}
where $\text{OFV}_0$ is the reference objective function value computed from the optimized parameter set. A large spread in $\Delta \text{OFV}$ across the perturbation range indicates a sensitive parameter, while a flat response indicates insensitivity. To rank \acp{mp} by their overall influence, the maximum absolute $\Delta \text{OFV}$ observed within the full perturbation range is used as the sensitivity score.
